\chapter{Présentation du contexte}
{
	\section{La Société pour l'Informatique Industrielle}
	{
		\subsection{Acteur de la transformation numérique}
		{
			\par
			{
				SII, ou Société pour l'Informatique Industrielle, est une entreprise de services du numérique française implantée dans dix-huit pays (voir \autoref{fig:sii_impl}) au travers de soixante-six agences. Elle propose de multiples services en \gls{b2b}, dont principalement l'étude et le conseil en technologies sur la maîtrise d'ouvrage de projets et la conception et l'intégration de systèmes.
			}
			
			\begin{figure}[h]
			{
				\centering
				\includegraphics[page=1,width=0.92\textwidth]{figures/sii_impl.pdf}
				\caption{Implantation de SII à travers le monde\cite{sii_rf}}
				\label{fig:sii_impl}
			}
			\end{figure}
			
			\par
			{
				L'activité au sein de l'entreprise s'organise en quatre unités organisationnelles (\emph{Business Units}, ou \gls{bu}), dont les noms se rapportent à leurs activités principales. Chacune s'intéresse à un groupe de clients majeurs de ses domaines d'activités:
				\begin{description}
					\item[BU Telecom:] Bouygues Télécom, SFR Numéricable, Huawei, Nokia\dots
					\item[BU Aero Space Defense:] MBDA, Thales, Nexter, Airbus, Dassault\dots
					\item[BU Energie Transport Retail:] EDF, Engie, Areva, SUEZ Evt., RATP, SNCF, Air France, FNAC, La Poste, Carrefour\dots
					\item[BU Banque Assurance Maladie:] Banque de France, BNP Paribas, BPCE, Société Générale, Crédit Mutuel, La Banque Postale, AXA, GMF, Swiss Life\dots
				\end{description}
			}
			
			\par
			{
				La fiche signalétique suivante donne une vue d'ensemble de l'entreprise et un contexte de réalisation du stage:
			
				\begin{description}
					\item[Nom de l'entreprise:] Société pour l'informatique industrielle (SII)
					\item[Date de création:] 1979\cite{sii_rf}
					\item[Statut juridique:] Société anonyme à directoire et conseil de surveillance
					\item[Siège social:] Paris
					\item[Effectif:] 5949\cite{sii_rs_2017}
					\item[Chiffre d'affaire:] 438,85M\euro\cite{sii_rs_2017}
					\item[Capital:] 103M\euro\cite{sii_rf}
					\item[Direction:] Patrice Demay (directeur général France), Eric Matteucci (président du directoire) et Jean-Paul Chevée (directeur général international)
					\item[Secteurs d'activité:] Aérospatiale, Banque, Défense, Énergie, Finance, Industrie, Média, Télécoms, Transports
					\item[Site internet:] \url{http://www.groupe-sii.com}
				\end{description}
			}
			
			\par
			{
				Le stage décrit ici s'est déroulé dans l'agence SII de Bourges, qui est rattachée à l'agence d'Ile-de-France. Elle compte une cinquantaine de collaborateurs affectés aux \gls{bu} \emph{Energie Transport Retail} et \emph{Aero Space Defense}, qui concentre la majorité des activités de l'agence, ainsi qu'au service Qualité, chargé de l'application des normes et du suivi de la conduite de projets.
			}
			
		}
		\subsection{Historique}
			
			\par
			{
				Fondée en \oldstylenums{1979}, SII a su se développer à l'international jusqu'à devenir un des leaders des SSII françaises. L'historique suivant permet d'appréhender sa progression au fil du temps:
			}
			
			\begin{description}
			{
				\item[1979:] Création de l'entreprise à Paris par Bernard Huvé, alors ingénieur consultant.
				\item[1984:] À l'occasion d'un contrat avec le laboratoire d'IBM, une deuxième agence est ouverte à Nice.
				\item[1987-1989:] Déploiement de SII en Île-de-France avec la création des deux agences de Cergy Pontoise et Vélizy.
				\item[1992:] À la suite d'évolutions importantes du secteur de l'informatique, l'entreprise se recentre sur l'assistance technique et obtient la certification ISO 9001.
				\item[1997:] Ouverture des agences de Rennes et Aix-en-Provence.
				\item[1999:] Introduction en bourse de l'entreprise.
				\item[2000-2007:] Accroissement du maillage national: ouverture d'agences à Bordeaux, Brest, Caen, La Ciotat, Lannion, Le Mans, Lyon, Niort, Montpellier, Tours et Vitrolles. Ouverture de filiales internationales en République Tchèque et en Pologne.
				\item[2010:] Acquisition de AIDA Development, société allemande, qui permet à l'entreprise de pénétrer le premier marché européen.
				\item[Août 2010:] Création de l'agence SII Bourges.
				\item[2015:] Implantation aux Pays-Bas et en Colombie.
				\item[2016:] Implantation aux Canada et en Angleterre.
			}
			\end{description}
	}
	\section{Développement et exposition du savoir-faire}
	
		\subsection{SII Research et le pôle robotique}
			
			\par
			{
				L'entreprise SII dispose d'un pôle innovation, incarné par le laboratoire \emph{SII Research}, qui vise à regrouper et structurer les projets relatifs à l'innovation technologique. Contrairement au fonctionnement des \gls{bu}, elle ne s'addresse pas à un secteur d'activité ou un client particulier, mais vise à développer les connaissances et le savoir-faire du groupe en effectuant des activités telles que de la veille technologique, la gestion de la propriété intellectuelle et la gestion de projets \gls{rd} internes. A l'occasion de la réalisation de ces derniers, les acteurs du pôle sont amenés à réaliser des \gls{poc}, activité dans laquelle s'inscrit le projet \gls{robo}. Parmi leurs projets achevés, on peut citer \emph{CATED}, une application d'aide à l'autonomisation pour personnes autistes, \emph{Entrepôt Augmenté}, un projet utilisant la réalité augmentée pour optimiser le temps de traitement des commandes et \emph{TryAplane}, un système d'aide à la formation aéronautique.
				
				\begin{figure}[h]
				{
					\centering
					\includegraphics[page=1,width=0.6\textwidth]{figures/siiresearch.png}
					\caption{Logo de la cellule SII Research}
					\label{fig:sii_research}
				}
				\end{figure}
			}
			\par
			{
				Les activités effectuées au long du projet \gls{robo} sont rattachées à une entitée inclue dans la cellule SII Research, le pôle robotique. Ce pôle est issu d'un partenariat entre l'agence SII Bourges et l'INSA Centre Val de Loire conclu en \oldstylenums{2015}. Son activité a déjà fait l'objet d'un stage en \oldstylenums{2016}, qui a consisté à réaliser un système de reconstruction d'environnement en 3D à partir d'images ou de flux vidéo\cite{qr3d}, destiné à un système robotique complet.
			}
			
		\subsection{Technologies de pointe}
		
			\par
			{
				En dehors de promouvoir l'entreprise et l'agence de Bourges auprès des étudiants de l'INSA Centre Val de Loire, ce partenariat permet de développer les connaissances et le savoir faire du groupe en mettant en relation proche le monde de la recherche scientifique avec celui de l'entreprise. En effet, ce projet met en jeu des techniques complexes et relativement récentes dans le monde de l'informatique et du traitement du signal. Sa réalisation a nécessité d'une part la mise en pratique de notions techniques avancées telles que la programmation \gls{gpu} et la présentation de scènes tridimensionnelles interactives en temps réel et d'autre part la compréhension et l'application d'algorithmes issus de l'apprentissage machine pour le \gls{slam} et la détection/classification.
			}
			\par
			{
				L'apprentissage profond, abordé dans ce rapport, est une technologie en plein essor qui soulève un engouement médiatique avéré et intéresse autant les chercheurs que les entreprises. En effet, cette discipline commence depuis peu à être intégrée dans des solutions informatiques professionnelles, notamment en matière de traduction automatique\cite{dltra}, de reconnaissance vocale\cite{dlrec}, de médecine\cite{dlmed}, de véhicules autonomes\cite{dlaut}\ldots \newline Les champs applicatifs paraissent intarissables, et la capacité de résolution de problèmes est telle qu'il semble de nos jours indispensable pour une entreprise de s'intéresser à ce savoir-faire.
			}
}
